\section{Introduction}\label{section:introduction}

Openly accessible data has long been shared through the Web; in
particular, many HTML pages contain \emph{entity-centric tables}, where
each line is about an entity, e.g., an album, and describes its attributes such as artist or year. Entity-centric tables have been leveraged to extract large-scale, open knowledge bases \cite{DBLP:conf/wsdm/LockardSDH20,DBLP:conf/semweb/AuerBKLCI07,DBLP:conf/esws/TanonWS20,DBLP:journals/ftdb/WeikumDRS21}, and for question answering \cite{DBLP:journals/pacmmod/WangF23,herzig-etal-2021-open,Christmann2023CompMixAB,Hulsebos2024ItTL}.
They are also frequent in \emph{data lakes}, \cite{fan_semantics-aware_2023,deng_lakebench_2024,christensen_fantastic_2025}.
In this work, we focus on \emph{openly accessible} Web data, independently of their licensing status. Among them, \emph{statistics datasets} (\textbf{SDs}, in short) form a distinct and highly valuable category.
They are compiled by domain specialists typically working for national
or international governing bodies, such as the
\href{https://www.un.org/}{United Nations}, the
\href{https://www.imf.org/}{IMF}, but also by companies, NGOs, etc.
Some organizations, e.g.,
\href{https://ec.europa.eu/eurostat/}{Eurostat}, publish hundreds of
thousands of SDs, updated yearly; others, e.g., an NGO focused on equitable justice, or a pollution watchdog in a certain area, may only publish a few dozen highly specialized datasets.

SDs~significantly differ from entity-oriented Web tables.
SD~content is
mostly numeric, e.g.,  birth and death counts by age in a
country over decades, or chip production per country and type of chip.
From a data management perspective, SDs are multidimensional
	aggregates, sometimes data cubes.
Also, SDs are mostly  published as standalone files, encoded in a
variety of formats, e.g., CSV, TSV, spreadsheets, XML dialects such as
SDMX~\cite{sdmx}, JSON, etc. Open SDs are also sometimes embedded in PDF
documents.

\emph{Answering queries or questions over tables}, and
  in our context, over SDs, is as the very core of the data management
  science and industry. The wide availability of statistic datasets on the Web requires tools that make
such data reachable in order to realize its potential; as the paper shows, such tools are currently lacking.

SDs are of great public utility. Officials rely on them to
inform decisions and support debates; social scientists use them to
analyze policy impacts and detect trends; newsrooms use them to check the
accuracy of public statements, make comparisons across countries, etc.
Also, the \emph{statistical literacy} of the general public needs to be improved: for instance, US voters vastly misrepresent the size of various fractions of
the US population, e.g., how many are transgender (estimated: 21\%,
true: 0.6\%)~\cite{wrongstats}. Thus, \emph{building and using warehouses of
statistics data} is central. Research works seeking to automate the recognition and checking of claims leveraging on retrieved corpora of SDs include, e.g.,~\cite{DBLP:journals/pvldb/Karagiannis0PT20,DBLP:journals/debu/0002P21,DBLP:conf/webdb/CaoMT18,balalau2024star, claimbuster} in the Information Retrieval and Data Management, and \cite{vlachos-riedel-2015-identification} in NLP.

Finding all SDs published by an organization is desirable in
order to study them as a whole, compare among organizations, populate a data lake for fact-checking, etc. For example, a work historian may want to retrieve all SDs published by the \href{https://www.ilo.org}{ILO}; a journalist may need all \href{https://ec.europa.eu/eurostat/}{Eurostat} SDs on poverty.  Unfortunately, current methods for this are either lacking, or quite inefficient.

Some websites publish an API to retrieve all their SDs, e.g.,
\href{https://ec.europa.eu/eurostat/}{Eurostat},
yet writing
custom code for each such API is cumbersome and brittle when APIs change.
Worse, many sites hosting numerous Open SDs, e.g., the ILO, do not
provide APIs to access all their data.
\emph{Search engines} (SEs, in short) and associated data
portals turn out to provide access to only a tiny fraction of
existing SDs (see Sec.~\ref{subsubsection:search_engines}); also, how these are selected is opaque to users.
A  \textbf{na\"ive, exhaustive website crawl is extremely inefficient} for large websites:
\begin{inparaenum}[(i)]
	\item acquiring a full
	website takes space and time;
	\item {\em crawling ethics} requires to wait
	(typically 1~second) between two successive HTTP requests; for
	a site of 1~million pages, such waits, alone, take 11
	days.
\end{inparaenum}
Thus, we are interested in a focused crawl, seeking to acquire within a given website as many \textbf{targets} as possible, while
\textbf{minimizing the resources consumed}, e.g., HTTP requests or
volume of transferred data.
Motivated by our SD retrieval problem, we define \emph{targets} as \emph{data files (CSV, spreadsheet, etc.)}. This can generalize to \emph{any} definition of target. Also, we aim for an efficient method, feasible on a single machine, as opposed to SE-scale parallel infrastructure.

Focused crawlers have been extensively
studied~\cite{chakrabarti1999focused,diligenti2000focused,gouriten2014scalable,meusel2014focused,faheem2015adaptive,han2018focused,kontogiannis2022tree}.
However, existing approaches often rely on \emph{assumptions that may not
hold in our setting}, e.g., that all pages on a
topic tend to be close within a website~\cite{10.1145/345508.345597}.
Early work includes generic focused crawling
approaches~\cite{chakrabarti1999focused,diligenti2000focused} (that we
will illustrate with a \textsf{\small FOCUSED} baseline); other study
large-scale website selection~\cite{meusel2014focused} (as opposed to SDs
within a given website), topic-driven page retrieval methods such as
\textsf{\small TRES}~\cite{kontogiannis2022tree}
or~\cite{chakrabarti1999focused,gouriten2014scalable,han2018focused}, and
non-topical objectives such as textual
diversity~\cite{faheem2015adaptive} (adapted to our setting
as a \textsf{\small TP-OFF} baseline).

Motivated by the large numbers of open, valuable SDs accessible
via navigation, we address the challenging problem of
\textbf{retrieving SDs in a website  as efficiently as possible, 
without relying on prior knowledge of the website's structure or content.}

\setlength{\fboxsep}{2pt}
\paragraph*{Contributions} Our contributions are as follows. \begin{inparaenum}[(1)]
	\noindent\item We formalize
	our graph crawling problem and show that optimally solving it is
	intractable (Sec.~\ref{sec:problem-model}).
	\noindent\item We propose a novel approach \textsf{\small SB-CLASSIFIER} based on
		reinforcement learning (RL) in Sec.~\ref{section:solution}, relying on two crucial  hypotheses: \begin{inparaenum}[(i)]
		\item 
		\emph{links found on similar \fbox{tag paths} within HTML web pages} (a tag path~\cite{miao2009extracting, furche2012turn} goes from the page root, to a hyperlink tag such as \texttt{\small $<$a$>$}) \emph{lead to similar content},
		\item \emph{we can learn, from the tag path structure, which
			tag paths are likely to lead to links towards targets}.
	\end{inparaenum}
	\noindent\item We demonstrate the effectiveness and efficiency of our approach
	through extensive experiments on diverse
		websites, totalling 22.2 millions pages (Sec.~\ref{sec:results}). \textsf{\small SB-CLASSIFIER} outperforms all baselines.
	On some websites we experimented on, in particular very large ones,
	\textbf{our crawler retrieves 90\% of the targets  accessing only 20\% of the
		webpages} (Sec.~\ref{subsection:dataset_characteristics}).
\end{inparaenum}

We discuss related work in
Sec.~\ref{section:related_work}. This paper is an extended version of the conference paper~\cite{gauquier2026efficient2}. A Github repository~\cite{github_repository} provides code to reproduce the experiments. 
